% !TeX root = ../教材题库_全部条目_转写工作稿.tex
% Source: XJ-B1-C05-EX-P179-Y0233
% 第5章 / 5.3 三角函数的图象与性质 / 教材 p.179 / PDF 第186页
% OCR未校对：公式与文字均可能有误；未完成数学TeX转写；禁止直接发给学生
% 原题及相邻原解截图（核对用）：../题目原图/XJ-B1-C05-EX-P179-Y0233-01.webp
% 本文件只在转写完成后作为题干使用；OCR 条目尚不是可靠题干。
\begin{bankitem}{5.3 例题 例题1}{第5章 / 5.3 三角函数的图象与性质 / 教材 p.179 / PDF 第186页}{XJ-B1-C05-EX-P179-Y0233}{OCR未校对：公式与文字均可能有误；未完成数学TeX转写；禁止直接发给学生}
\begin{ocrdraft}
[教材p.179 / PDF第186页]
\par
0]用“五点法”画出下列函数的简图：
\par
出余弦曲线，一cos在
\par
（1）訓一1十sin心e[0，2丌]；
\par
[0，2丌]上的关键五点吗？
\par
（2）y=2cos，工e[0，2冠
\par
解（D按五个关键点列表：
\par
3丌
\par
0
\par
2
\par
2
\par
0
\par
1
\par
0
\par
0
\par
Sln工
\par
1+sin工
\par
1
\par
2
\par
1
\par
0
\par
1
\par
描点，并将这些点依次连成一条光滑曲线，即得所求图
\par
象，如图5，3一5．
\par
你能从甬数图象变换
\par
的角度，利用函数，一
\par
y=l+sinx
\par
sin，工e[0，2丌]的图象
\par
来得到3，一1+sin宀xe
\par
[0，2丌]的图象吗？
\par
0
\par
2是
\par
2
\par
=S1nx
\par
图5．3一5
\par
（2）按五个关键点列表：
\par
3丌
\par
0
\par
2
\par
2
\par
COS工
\par
1
\par
0
\par
0
\par
1
\par
2C05工
\par
2
\par
2
\par
0
\par
0
\par
描点，并将这些点依次连成一条光滑曲线，即得所求图象，如图5．3
\par
一6．
\par
y=2cosX
\par
2
\par
0
\par
2兀
\par
2
\par
2
\par
图5．3一6
\end{ocrdraft}
\end{bankitem}
